\documentclass[../main.tex]{subfiles}
\begin{document}

\section{Introduction}
\section{Attacks}
\subsection{White Box Attacks}
\subsubsection{FGSM Attack on UNet}

\paragraph{Introduction}\leavevmode\\
The Fast Gradient Sign Method (FGSM) attack is applied to UNet model to evaluate its robustness by degrading segmentation performance. FGSM introduces $\ell_\infty$-bounded perturbations, controlled by $\epsilon$, to input images, causing reduced accuracy and Dice scores compared to ground truth masks, often resulting in incorrect segmentations or missed regions. This attack preserves visual similarity to the original image while disrupting the model’s predictions, highlighting vulnerabilities in critical applications like medical imaging.

\paragraph{Methodology}\leavevmode\\

\begin{breakablealgorithm}{Fast Gradient Sign Method (FGSM) Attack for UNet}
\begin{algorithmic}[1]
\Require $model$: Trained UNet model, $Data$: (image, mask) pairs, $\epsilon$: Perturbation magnitude
\Ensure $adversarial\_image$: Adversarial image batch, normalized to [0, 1]

\State Initialize loss function: $loss\_fn \gets \text{BCEWithLogitsLoss}()$
\State Set $model$ to evaluation mode

\State Initialize $adversarial\_image \gets image.clone().detach()$
\State Enable gradient computation: $adversarial\_image.requires\_grad \gets \text{True}$

\State Compute prediction: $prediction \gets model(adversarial\_image)$
\State Compute loss: $loss \gets loss\_fn(prediction, mask.unsqueeze(1))$
\State Backpropagate: $loss.backward()$

\State Extract gradient: $grad \gets adversarial\_image.grad$
\State Update image: $adversarial\_image \gets adversarial\_image + \epsilon \cdot \text{sign}(grad)$
\State Clip to valid range: $adversarial\_image \gets \text{clamp}(adversarial\_image, 0, 1)$

\State Detach $adversarial\_image$ from computation graph
\State \Return $adversarial\_image$
\end{algorithmic}
\end{breakablealgorithm}
\paragraph{Results}\leavevmode\\

The FGSM attack was applied with perturbation magnitudes $\epsilon=0.031$ and $0.063$. Figure~\ref{fig:fgsm_example_unet} illustrates the effect of the attack on image example from the caravana dataset, showing the original image, the adversarial image for \epsilon=0.031$ and $0.063$.

% Figure: FGSM Example Image

\begin{figure}[H]
    \centering
    \begin{subfigure}[b]{0.2\textwidth}
        \includegraphics[width=\textwidth]{images/CV/Segmentation/original_unet.jpg}
        \caption{Original Input Image}
    \end{subfigure}
    \hfill
    \begin{subfigure}[b]{0.2\textwidth}
        \includegraphics[width=\textwidth]{images/CV/Segmentation/fgsm_unet_eps=0.031.jpg}
        \caption{Adversarial image. $\epsilon$ = 0.031}
    \end{subfigure}
        \hfill
    \begin{subfigure}[b]{0.2\textwidth}
        \includegraphics[width=\textwidth]{images/CV/Segmentation/fgsm_unet_eps=0.062.jpg}
        \caption[si]{Adversarial image. $\epsilon$ = 0.062}
    \end{subfigure}

    \vspace{0.2cm} % Adding space between rows

    \begin{subfigure}[b]{0.2\textwidth}
        \includegraphics[width=\textwidth]{images/CV/Segmentation/mask_unet.jpg}
        \caption{Ground Truth Mask}
    \end{subfigure}
    \hfill
    \begin{subfigure}[b]{0.2\textwidth}
        \includegraphics[width=\textwidth]{images/CV/Segmentation/pred_fgsm_unet_eps=0.031.jpg}
        \caption{Prediction. $\epsilon$ = 0.031}
    \end{subfigure}
    \hfill
    \begin{subfigure}[b]{0.2\textwidth}
        \includegraphics[width=\textwidth]{images/CV/Segmentation/pred_fgsm_unet_eps=0.062.jpg}
        \caption{Prediction. $\epsilon$ = 0.062}
    \end{subfigure}
    \label{fgsm_example_unet}
\end{figure}


The following table ~\ref{tab:unet_fgsm_overall_performance} presents the overall performance metrics (e.g., accuracy and Dice score) on the Caravana validation set, highlighting the impact of FGSM attack on detection performance.



% Table 1: Overall Performance
\begin{table}[H]
\caption{Overall Performance of UNet under Clean and FGSM Attacks}
    \label{tab:unet_fgsm_overall_performance}
    \centering
    \footnotesize
    \begin{tabular}{@{\hspace{0.3cm}}lS[table-format=1.3]S[table-format=1.3]S[table-format=1.3]S[table-format=1.3]@{\hspace{0.3cm}}}
        \toprule
        \textbf{Metric} & {\textbf{Clean}} & {\textbf{FGSM $\epsilon=0.031$}} & {\textbf{FGSM $\epsilon=0.062$}}  \\
        \midrule
        Accuracy (\%)    & 99.31 & 39.55 & 24.89 \\
        Dice Score        & 0.98 & 0.41 & 0.36  \\
        \bottomrule
    \end{tabular}
\end{table}
Under clean conditions, the model achieves a high accuracy of 99.31\% and a Dice score of 0.98, reflecting excellent detection performance. With the FGSM attack at $\epsilon$ = 0.031, accuracy drops to 39.55\% and the Dice score decreases to 0.406, showing a notable degradation. The effect is even more pronounced at $\epsilon$ = 0.062, where accuracy falls to 24.89\% and the Dice score drops to 0.36, indicating severe disruption in detection capability. The visual comparison in Figure~\ref{fig:fgsm_example_unet} further supports this, with increasing ε values leading to poorer prediction masks, highlighting the vulnerability of the model to adversarial perturbations.

\subsubsection{PGD Attack on UNet}

\paragraph{Introduction}\leavevmode\\
The Projected Gradient Descent (PGD) \cite{madry2018towardspgd} attack is applied to the UNet model to evaluate its robustness by degrading segmentation performance. Unlike the single-step Fast Gradient Sign Method (FGSM), which applies a one-time perturbation, PGD iteratively refines perturbations using gradient updates and projects them onto an $\ell_\infty$-bounded ball, achieving stronger disruptions. This results in reduced accuracy and Dice scores compared to ground truth masks, often causing incorrect segmentations or missed regions. PGD preserves visual similarity to the original image while significantly disrupting the model’s predictions.

\paragraph{Methodology}\leavevmode\\

\begin{breakablealgorithm}{Projected Gradient Descent (PGD) Attack for UNet}
\begin{algorithmic}[1]
\Require $model$: Trained UNet model, $data$: Data loader yielding (image, mask) pairs, $\epsilon$: Perturbation magnitude, $\alpha$: Step size, $iterations$: Number of iterations
\Ensure $adversarial\_image$: Adversarial image batch, normalized to [0, 1]

\State Initialize loss function: $loss\_fn \gets \text{BCEWithLogitsLoss}()$
\State Set $model$ to evaluation mode

\State Initialize $adversarial\_image \gets image.clone().detach()$
\For{$i \gets 1$ to $iterations$}
    \State Compute prediction: $prediction \gets model(adversarial\_image)$
    \State Compute loss: $loss \gets loss\_fn(prediction, mask)$
    \State Backpropagate: $loss.backward()$
    
    \State Update image: $adversarial\_image \gets adversarial\_image + \alpha \cdot \text{sign}(adversarial\_image.grad)$
    \State Project to $\epsilon$-ball: $adversarial\_image \gets \max(\min(adversarial\_image, original\_image + \epsilon), original\_image - \epsilon)$
    \State Clip to valid range: $adversarial\_image \gets \text{clamp}(adversarial\_image, 0, 1)$
\EndFor

\State \Return $adversarial\_image$

\end{algorithmic}
\end{breakablealgorithm}
\paragraph{Results}\leavevmode\\

The PGD attack was applied with the maximum allowable perturbation  $\epsilon$=0.031 and 0.062 with the step size used for each gradient update $\alpha$=2/255 for 10 iterations. Figure~\ref{fig:pgd_example_unet} illustrates the effect of the attack on image example from the caravana dataset, showing the original image, the adversarial image for $\epsilon$=0.031 and 0.062.

% Figure: FGSM Example Image

\begin{figure}[H]
    \centering
    \begin{subfigure}[b]{0.2\textwidth}
        \includegraphics[width=\textwidth]{images/CV/Segmentation/original_unet.jpg}
        \caption{Original Input Image}
    \end{subfigure}
    \hfill
    \begin{subfigure}[b]{0.2\textwidth}
        \includegraphics[width=\textwidth]{images/CV/Segmentation/pgd_unet_eps=0.031.jpg}
        \caption{Adversarial image. $\epsilon$ = 0.031, $\alpha$ = 2/255}
    \end{subfigure}
        \hfill
    \begin{subfigure}[b]{0.2\textwidth}
        \includegraphics[width=\textwidth]{images/CV/Segmentation/pgd_unet_eps=0.062.jpg}
        \caption[si]{Adversarial image. $\epsilon$ = 0.062, $\alpha$ = 2/255}
    \end{subfigure}

    \vspace{0.2cm} % Adding space between rows

    \begin{subfigure}[b]{0.2\textwidth}
        \includegraphics[width=\textwidth]{images/CV/Segmentation/mask_unet.jpg}
        \caption{Ground Truth Mask}
    \end{subfigure}
    \hfill
    \begin{subfigure}[b]{0.2\textwidth}
        \includegraphics[width=\textwidth]{images/CV/Segmentation/pred_pgd_unet_eps=0.031.jpg}
        \caption{Prediction. $\epsilon$ = 0.031, $\alpha$ = 2/255}
    \end{subfigure}
    \hfill
    \begin{subfigure}[b]{0.2\textwidth}
        \includegraphics[width=\textwidth]{images/CV/Segmentation/pred_pgd_unet_eps=0.062.jpg}
        \caption{Prediction. $\epsilon$ = 0.062, $\alpha$ = 2/255}
    \end{subfigure}
    \label{pgd_example_unet}
\end{figure}


The following table ~\ref{tab:unet_fgsm_overall_performance} presents the overall performance metrics (e.g., accuracy and Dice score) on the Caravana validation set for  $\epsilon$=0.031 and 0.062 with $\alpha$=2/255 for 10 iterations, highlighting the impact of PGD attack on detection performance.



% Table 1: Overall Performance
\begin{table}[H]
\caption{Overall Performance of UNet under Clean and PGD Attacks}
    \label{tab:unet_fgsm_overall_performance}
    \centering
    \footnotesize
    \begin{tabular}{@{\hspace{0.3cm}}lS[table-format=1.3]S[table-format=1.3]S[table-format=1.3]S[table-format=1.3]@{\hspace{0.3cm}}}
        \toprule
        \textbf{Metric} & {\textbf{Clean}} & {\textbf{PGD $\epsilon=0.031$}} & {\textbf{PGD $\epsilon=0.062$}}  \\
        \midrule
        Accuracy (\%)    & 99.31 & 37.21 & 30.32 \\
        Dice Score        & 0.98 & 0.396 & 0.371  \\
        \bottomrule
    \end{tabular}
\end{table}
The PGD attack results on the Caravana validation set, as shown in Table 2, reveal a significant decline in the UNet model's performance compared to clean conditions, with accuracy dropping from 99.31\% to 37.21\% at $\epsilon$ = 0.031 ($\alpha$ = 2/255) and further to 30.32\% at $\epsilon$ = 0.062 ($\alpha$ = 2/255) over 10 iterations, alongside Dice scores falling from 0.98 to 0.396 and 0.371, respectively. This indicates a severe disruption in detection capability, with the visual comparison in Figure~\ref{fig:pgd_example_unet} showing poorer prediction masks as $\epsilon$ increases, underscoring the model's vulnerability to adversarial perturbations. Notably, the perturbations in PGD are less visible compared to FGSM, yet they still cause substantial degradation.

\subsubsection{MI-FGSM Attack on UNet}

\paragraph{Introduction}\leavevmode\\\subsection{Black Box Attacks}
\section{Defenses}


\end{document}
